\section{Introduction}

Modern machine learning systems are typically designed to produce a single output given an input query. Formally, this corresponds to selecting a representative element from a conditional distribution (P(y \mid q)), often through maximum likelihood or averaging. While effective in many settings, this paradigm implicitly assumes that the underlying semantic space is well-represented by a single point.

However, in many real-world scenarios, this assumption fails. When multiple valid but conflicting interpretations exist—such as in opinionated text, multi-agent reasoning, or cross-cultural data—collapsing the distribution into a single output discards essential structural information. In particular, relationships between alternative outputs, including agreement, contradiction, and partial overlap, are lost.

Recent work has attempted to address this limitation by modeling distributions or sets of outputs. Nevertheless, these approaches largely operate at the level of probability or geometry, without capturing the higher-order structure that governs how local semantic relationships combine into global consistency.

In this work, we propose a fundamentally different perspective: we represent meaning not as a point or distribution, but as a \emph{topological invariant}. Specifically, we introduce \emph{persistent sheaf cohomology}, a framework that integrates topological data analysis (TDA) and sheaf theory to capture both geometric structure and semantic consistency across multiple scales.

Given a collection of candidate outputs, we first embed them into a metric space and construct a filtration of simplicial complexes. We then equip each complex with a sheaf that encodes local semantic relationships, allowing us to model how meanings interact across overlapping contexts. By computing cohomology groups across the filtration, we obtain a persistence module that captures stable structural features.

A key insight of our framework is that the first cohomology group (H^1) provides a natural measure of semantic inconsistency. Intuitively, non-zero (H^1) indicates the presence of cycles that cannot be resolved into a globally consistent interpretation, corresponding to conflicting meanings that cannot be reconciled. This interpretation allows us to move beyond heuristic measures of disagreement and instead define inconsistency as a mathematically well-founded invariant.

We support this perspective both theoretically and empirically. We prove that persistent sheaf cohomology is stable under perturbations of the embedding space, ensuring robustness to noise. Furthermore, we show that under mild conditions, the vanishing of (H^1) is equivalent to the existence of a globally consistent semantic assignment. Empirically, we demonstrate that the magnitude of (H^1) strongly correlates with human-annotated disagreement across multiple datasets, while conventional averaging-based methods fail to capture such structure.

Our results suggest that meaning is inherently structured and cannot be faithfully represented by a single value. By treating semantic inference as the extraction of topological invariants, we provide a new foundation for modeling multi-perspective data, with potential applications in uncertainty-aware reasoning, multi-agent systems, and interpretability.

\paragraph{Contributions.}
We summarize our main contributions as follows:
(i) we introduce persistent sheaf cohomology as a unified framework for semantic inference;
(ii) we provide theoretical guarantees, including stability and a cohomological characterization of inconsistency;
(iii) we empirically demonstrate that our method captures semantic structure that is lost in standard approaches.
